\clearpage
\phantomsection% 
\addcontentsline{toc}{chapter}{RINGKASAN}
\thispagestyle{fancy}

\begin{center}
	\large \bfseries \MakeUppercase{Ringkasan}\\
	\normalsize \normalfont {\thetitle}\\
	\normalsize \normalfont {\theauthor}\\
	\bigskip
	
	\normalsize \normalfont \justifying
	\normalsize \normalfont \justifying
	Kemajuan teknologi astronomi telah mendorong penggunaan teleskop robotik yang mampu melacak objek langit secara otomatis dengan akurasi tinggi. Namun, teleskop otomatis yang tersedia di pasaran sering menghadapi tantangan pada tahap konfigurasi awal, khususnya dalam proses penyelarasan (\textit{alignment}) dan kalibrasi yang memakan waktu serta memerlukan keterampilan teknis. Permasalahan ini membatasi partisipasi masyarakat, terutama pelajar dan komunitas astronomi amatir, dalam pengembangan ilmu astronomi. Oleh karena itu, penelitian ini menawarkan solusi berupa pengembangan teleskop robotik berbasis teknologi lokal dengan dudukan \textit{alt-azimuth} yang dilengkapi sistem kontrol otomatis melalui aplikasi Android, sensor IMU MPU6050, magnetometer QMC5883L, dan encoder magnetik AS5600 untuk estimasi orientasi secara \textit{real-time}.\par
	
	Penelitian ini difokuskan pada dua permasalahan utama, yaitu bagaimana merancang dan membangun teleskop robotik dua sumbu berbasis \textit{Alt-Azimuth} yang dapat menyesuaikan posisi akibat pergeseran atau perpindahan serta mendukung pelacakan objek antariksa, dan bagaimana hasil evaluasi integrasi sensor IMU dengan penerapan \textit{Mahony Filter} dalam meningkatkan akurasi estimasi sudut pada teleskop robotik. Tujuan penelitian adalah merancang dan membangun teleskop robotik dua sumbu berbasis \textit{Alt-Azimuth} yang mampu menyesuaikan posisi saat terjadi pergeseran atau perpindahan serta dapat melakukan pelacakan terhadap objek antariksa, dan mengevaluasi kinerja integrasi sensor IMU dengan penerapan \textit{Mahony Filter} dalam meningkatkan akurasi estimasi sudut pada sistem teleskop robotik.\par
	
	Penelitian dilaksanakan melalui beberapa tahapan sistematis dimulai dengan studi literatur komprehensif mengenai teori teleskop reflektor, sistem koordinat langit, dan algoritma \textit{Mahony Filter}. Observasi lapangan dan wawancara dilakukan di Observatorium Astronomi ITERA Lampung (OAIL) untuk memperoleh informasi kebutuhan pengguna dan kendala operasional. Pengembangan prototipe mengikuti model \textit{Prototyping} iteratif yang mencakup analisis kebutuhan, desain sistem mekanik, perancangan rangkaian elektronik berbasis ESP32 dengan arsitektur \textit{dual-core}, pengembangan perangkat lunak, dan implementasi algoritma estimasi orientasi. Sistem menggunakan sensor IMU MPU6050 untuk pembacaan giroskop dan akselerometer, magnetometer QMC5883L untuk kompensasi \textit{heading}, dan encoder AS5600 untuk posisi absolut. Data sensor diproses menggunakan \textit{Mahony Filter} berbasis pengendali PI (Proporsional-Integral) dengan parameter gain kp dan ki untuk mengatasi \textit{noise}, \textit{drift}, dan ketidakakuratan sensor. Evaluasi kinerja dilakukan menggunakan metrik standar deviasi untuk mengukur akurasi estimasi sudut, ketepatan \textit{pointing}, kemampuan \textit{tracking}, dan stabilitas sistem.\par
	
	Teleskop robotik berhasil dirancang dan dibangun dengan mengintegrasikan sensor IMU, magnetometer, dan encoder untuk memberikan informasi orientasi secara \textit{real-time}. Sistem menggunakan mikrokontroler ESP32 dengan arsitektur \textit{dual-core} yang memproses pembacaan sensor pada \textit{core} 0 dengan frekuensi 100 Hz dan kendali motor pada \textit{core} 1. Hasil evaluasi menunjukkan \textit{Mahony Filter} sangat efektif dalam meningkatkan akurasi estimasi sudut dengan mereduksi \textit{noise} sensor MPU6050 secara signifikan. Standar deviasi Roll menurun sebesar 79,1\% (dari 0,510$^{\circ}$ menjadi 0,107$^{\circ}$) dan Pitch menurun sebesar 84,2\% (dari 1,719$^{\circ}$ menjadi 0,271$^{\circ}$), sehingga menghasilkan estimasi orientasi yang stabil dan akurat. Sistem mampu menyesuaikan posisi secara otomatis saat terjadi pergeseran melalui mekanisme inisialisasi orientasi menggunakan pembacaan \textit{heading} dari magnetometer dan pembacaan sudut absolut dari encoder. Sistem pengarahan otomatis (\textit{GoTo}) berhasil mengarahkan teleskop ke objek langit dengan error berkisar 0$^{\circ}$10' hingga 2$^{\circ}$30'. Algoritma astronomi untuk perhitungan posisi Matahari, Bulan, dan Planet menghasilkan akurasi dengan error $< 0,3^{\circ}$. Kemampuan pelacakan (\textit{tracking}) objek antariksa telah berhasil diimplementasikan dengan sistem mampu mengikuti pergerakan semu benda langit. Pengujian pada empat objek dengan durasi pengamatan yang bervariasi menunjukkan bahwa sistem dapat menghitung dan mengkompensasi pergerakan koordinat horizontal secara kontinu. Pengamatan Matahari selama 19 menit menunjukkan objek masih dalam bidang pandang, namun pelacakan objek mendekati zenith seperti Bulan mengalami drift signifikan yang menyebabkan objek keluar dari medan pandang.\par
	
	Penelitian ini telah mengembangkan teleskop robotik berbasis teknologi lokal menggunakan sistem \textit{Alt-Azimuth} dengan integrasi sensor fusion dan \textit{Mahony Filter} yang mampu melakukan pengarahan objek langit secara otomatis. \textit{Mahony Filter} terbukti sangat efektif dalam mengatasi kelemahan sensor MPU6050 dengan menyediakan data orientasi yang stabil untuk sistem kendali motor dan transformasi koordinat. Sistem berhasil mengatasi keterbatasan konfigurasi awal teleskop konvensional melalui mekanisme kalibrasi yang lebih sederhana dan kemampuan adaptasi terhadap pergeseran posisi. Sistem pengarahan dapat mengarahkan teleskop dengan error berkisar 0$^{\circ}$10' hingga 2$^{\circ}$30' yang memadai untuk objek dengan diameter sudut besar. Untuk pelacakan objek, sistem mampu menghitung dan mengkompensasi pergerakan semu benda langit secara kontinu dengan keberhasilan yang bervariasi tergantung pada posisi objek di langit. Keterbatasan pelacakan disebabkan oleh kombinasi faktor mekanik (\textit{flex} struktur PLA dan \textit{backlash} gear), frekuensi update terbatas (5-10 detik), dan akumulasi error. Untuk aplikasi astronomi amatir dan edukasi dengan objek langit terang berdiameter besar pada posisi tidak mendekati zenith, sistem sudah cukup memadai, namun masih diperlukan perbaikan mekanik signifikan untuk pengamatan objek titik seperti planet dan bintang yang memerlukan presisi lebih tinggi.\par
	
	Pengembangan lebih lanjut dapat dilakukan dengan mengganti material PLA cetak 3D dengan material yang lebih \textit{rigid} seperti aluminium untuk mengurangi \textit{flex} mekanik, meningkatkan frekuensi pembaruan posisi dari 5-10 detik menjadi 1-2 detik untuk mengurangi akumulasi error, mengintegrasikan kamera \textit{guide} dengan algoritma deteksi objek untuk koreksi posisi secara \textit{real-time}, mengembangkan sistem kalibrasi multi-titik berbasis interpolasi atau regresi polinomial, mengurangi \textit{backlash} mekanik dengan gear presisi tinggi atau mekanisme \textit{pre-load}, menerapkan kontroler PID (\textit{Proportional-Integral-Derivative}) untuk sistem kendali motor pada kedua sumbu Az dan Alt yang dapat mengurangi \textit{error steady-state}, mempercepat waktu \textit{settling}, dan meminimalkan \textit{overshoot} saat pengarahan dengan parameter yang di-\textit{tuning} secara spesifik untuk setiap sumbu mengingat karakteristik beban berbeda, serta mengimplementasikan model \textit{ephemeris} yang lebih akurat dengan koreksi perturbasi planet, nutasi, aberasi, dan koreksi atmosferik untuk meningkatkan presisi hingga level \textit{arcminute} atau \textit{arcsecond}.\par

	\vfill
	
\end{center}
\clearpage